% EWAT — Slides minimales défensives (Phase E du plan)
% À utiliser en complément du rapport principal. Compile avec :
%   pdflatex defense_slides.tex (2 passes pour TOC + références)

\documentclass[10pt,aspectratio=169]{beamer}

\usepackage[utf8]{inputenc}
\usepackage[T1]{fontenc}
\usepackage[french]{babel}
\usepackage{booktabs}
\usepackage[table]{xcolor}
\usepackage{colortbl}
\usepackage{tabularx}
\usepackage{amsmath,amssymb}

\definecolor{ewatred}{RGB}{200,40,40}
\definecolor{ewatblue}{RGB}{30,80,140}
\definecolor{ewatgreen}{RGB}{40,140,60}

\usetheme{Madrid}
\usecolortheme{seahorse}
\setbeamertemplate{footline}[frame number]
\setbeamertemplate{navigation symbols}{}

\title[EWAT — Défense]{EWAT : réponses anticipées aux critiques}
\subtitle{3 slides défensives pour la soutenance}
\author{Wassim Badraoui}
\institute{Devoteam --- EPITA}
\date{Mai 2026}

\begin{document}

\begin{frame}
  \titlepage
\end{frame}

% =============================================================================
% SLIDE 1 — Verdict synthétique des stress tests
% =============================================================================
\begin{frame}{Slide 1 --- Verdict synthétique : stress tests H3}
  \framesubtitle{Pourquoi l'AUROC=0.987 n'est pas la métrique défensive}

  \begin{block}{La critique soulevée}
    \emph{« L'AUROC=0.987 est trop bon et suspect. »}
  \end{block}

  \begin{block}{La réponse en 5 stress tests}
    \scriptsize
    \begin{tabularx}{\linewidth}{l X l}
      \toprule
      \textbf{Test} & \textbf{Question} & \textbf{Verdict} \\
      \midrule
      \textbf{A1} distant-window     & Précursion temporelle sur labels EWAT ?   & \textcolor{ewatred}{$\Delta = -0.007$ → fuite} \\
      \textbf{A2} LOSO precursor     & Généralisation à un type inédit ?         & \textcolor{ewatred}{top-1 = 0.51 polarisé} \\
      \textbf{A3} permutation        & Signal réel ou bruit ?                    & \textcolor{ewatgreen}{$p < 0.01$ → signal réel} \\
      \textbf{A4} filtre $n_\text{pos} \geq 5$ & Clusters statistiquement reportables ?   & 5/10 (AUROC moyen = 0.975) \\
      \textbf{A5} paired $\Delta(B_4 - B_3)$ & STGCN apporte-t-il du gain ?               & \textcolor{ewatred}{IC=[$-0.03$, $+0.04$] contient 0} \\
      \bottomrule
    \end{tabularx}
  \end{block}

  \begin{block}{Conclusion}
    L'AUROC=0.987 mesure la \textbf{cohérence interne} du clustering EWAT (cible
    auto-référente), pas la précursion temporelle. Cette critique est
    \textbf{reconnue et adressée} par les Phases B et C ci-après.
    \\[0.3em]
    Scripts : \texttt{experiments/h3\_robustness/}.
  \end{block}
\end{frame}

% =============================================================================
% SLIDE 2 — Headline honnête
% =============================================================================
\begin{frame}{Slide 2 --- Headline honnête : 0.92 sur cible Chaos Mesh}
  \framesubtitle{Évaluation indépendante avec IC bootstrap explicites}

  \begin{block}{Pivot méthodologique (Phase B + C)}
    Remplacer la cible (cluster EWAT) par les \textbf{15 scénarios Chaos Mesh}
    (vérité terrain indépendante du pipeline).
  \end{block}

  \begin{block}{Headline défensif final (multi-seed, 10 graines, v4\_strat)}
    \centering
    \begin{tabular}{lcl}
      \toprule
      \textbf{Évaluation} & \textbf{macro-AUROC} & \textbf{IC 95\,\% / std} \\
      \midrule
      H3 EWAT-labels (\textbf{circulaire})    & $0.990 \pm 0.012$ & 10 graines \\
      \rowcolor{black!8}
      \textbf{B2 LR-OvR Chaos Mesh stratified} & $\mathbf{0.9201}$ & \textbf{[0.878, 0.956]} \\
      \rowcolor{black!8}
      \textbf{B2 LOSO macro-AUROC}             & $\mathbf{0.9298}$ & 15 folds \\
      H1 silhouette test (10 graines)          & $0.691 \pm 0.115$ & range $[0.52, 0.84]$ \\
      A1 $\Delta$(far$-$near) EWAT             & $-0.012 \pm 0.022$ & LEAK 9/10 \\
      \bottomrule
    \end{tabular}
  \end{block}

  \begin{block}{Précursion temporelle confirmée (C2-A1)}
    Sur le modèle STGCN entraîné directement sur Chaos Mesh, \textbf{$\Delta($far\,$-$\,near$) = -0.116$}
    --- la dynamique pré-injection compte pour 12\,pp d'AUROC. La fuite vue par A1
    était un artefact de la circularité, pas du signal.
  \end{block}
\end{frame}

% =============================================================================
% SLIDE 3 — Reframe des contributions
% =============================================================================
\begin{frame}{Slide 3 --- Reframe des contributions}
  \framesubtitle{Quatre axes défendables (et un résultat négatif honnête)}

  \begin{enumerate}
    \item \textbf{Géométrique} (contribution principale et défendable)
      \begin{itemize}
        \item H1 silhouette = $0.782 \pm 0.065$ (10 graines, min $0.62$)
        \item Mismatch corrigé : average+cosine vs ward+euclidean → $+51\%$
      \end{itemize}

    \item \textbf{Ontologique} (Phase 8, OWL/RDF formelle)
      \begin{itemize}
        \item 29 classes ancrées littérature (Soldani \& Brogi, Fu et al., Gregg, Aniello)
        \item 143 individus ABox, 11 object properties, raisonneur HermiT cohérent en 0.61s
        \item 3 causales (BH-FDR), 19 co-occurrences, 46 edges propagation
      \end{itemize}

    \item \textbf{Méthodologique} (stress tests + open-set)
      \begin{itemize}
        \item 5 stress tests (A1--A5) qui transforment la critique en transparence
        \item OpenMax/EVT pour la détection de scénarios inédits (Bendale \& Boult 2016)
      \end{itemize}

    \item \textbf{Prédictive honnête} (Phase B/C)
      \begin{itemize}
        \item macro-AUROC $= \mathbf{0.920}$ sur Chaos Mesh v4\_strat, IC [0.878, 0.956]
        \item Précursion temporelle confirmée ($\Delta(\text{far}-\text{near})=-0.116$)
      \end{itemize}
  \end{enumerate}

  \vspace{0.4em}
  \begin{block}{Résultat négatif assumé}
    \scriptsize H2a look-through MMD$^2$ \textbf{définitivement falsifié}
    (v3 : $p = 0.27$, v4\_strat : $p = 0.37$). Mécanisme falsifié, pas juste la
    durée d'épisode. Résultat scientifique honnête.
  \end{block}
\end{frame}

% =============================================================================
% Anticipation des questions
% =============================================================================
\begin{frame}{Bonus --- Anticipation des questions}

  \begin{block}{« Votre AUROC=0.987 est suspect »}
    Oui, et c'est pourquoi nous reportons \textbf{0.920 sur Chaos Mesh indépendant
    (slide~2)}. Le 0.987 est explicitement étiqueté comme cohérence interne du clustering.
  \end{block}

  \begin{block}{« L'encodeur STGCN n'aide pas »}
    Vrai en agrégé (A5 paired $\Delta$ IC contient 0). Sa valeur est \textbf{géométrique}
    (H1 sil = 0.78) et \textbf{ontologique} (Phase 8). En prédiction agrégée, un
    LR sur features brutes instance-normalisées suffit (B2 = 0.920).
  \end{block}

  \begin{block}{« Vous évaluez sur les scénarios d'entraînement »}
    Oui pour H3, et c'est pourquoi nous ajoutons :
    \begin{itemize}
      \item LOSO macro-AUROC = 0.930 ± 0.007 (15 folds, retrain par fold)
      \item OpenMax/EVT pour les scénarios totalement inédits (Phase C3)
      \item LOSO top-1 = 0 reconnu comme limitation tautologique d'un classifieur fermé
    \end{itemize}
  \end{block}

  \begin{block}{« Pourquoi ewat\_v4\_strat et pas ewat\_v4 ? »}
    Le split temporel original ewat\_v4 avait 4 scénarios entièrement absents du
    training → AUROC=0.500 trivial. Le split stratifié-temporel v4\_strat (270/60/45)
    rétablit une évaluation honnête.
  \end{block}

\end{frame}

% =============================================================================
% SLIDE 5 — Validation opérationnelle : latence + power
% =============================================================================
\begin{frame}{Slide 5 --- Validation opérationnelle}
  \framesubtitle{Latence end-to-end (C-3) + power statistique (C-5)}

  \begin{columns}[T]
    \begin{column}{0.55\textwidth}
      \begin{block}{Latence E2E (200 itérations CPU)}
        \scriptsize
        \begin{tabular}{lrrr}
          \toprule
          \textbf{Étape} & \textbf{médiane} & \textbf{p95} & \textbf{budget} \\
          \midrule
          Drift (0) & 0.01 & 0.01 & 1000 ms \\
          Encoder+Siamois (1+2) & 0.96 & 1.97 & 2000 ms \\
          Précurseurs (3) & 1.85 & 3.91 & 1000 ms \\
          \rowcolor{black!8}
          \textbf{TOTAL} & $\mathbf{9.3}$ & $\mathbf{13.3}$ & $\mathbf{5000}$ \\
          \bottomrule
        \end{tabular}
        \vspace{0.3em}
        Verdict : \textcolor{ewatgreen}{\textbf{\textbullet{} 375$\times$ sous budget}}
      \end{block}
    \end{column}

    \begin{column}{0.45\textwidth}
      \begin{block}{Power statistique (Hanley-McNeil)}
        \scriptsize
        \begin{itemize}
          \item 5/10 clusters reportables ($n_\text{pos} \geq 5$)
          \item Power moyenne sur reportables : $1.0$
          \item C6, C9 ($n_\text{pos} = 1$) marqués \emph{non concluant}
        \end{itemize}
        \vspace{0.3em}
        \textbf{Cible future} : $n_\text{pos} \geq 17$ pour détecter AUROC $= 0.7$ avec power $> 0.8$.
      \end{block}
    \end{column}
  \end{columns}

  \begin{block}{Conclusion}
    Le pipeline est \textbf{opérationnellement viable} (SLA $< 5$\,s largement
    respecté) et statistiquement défensable sur les 5 clusters reportables.
    Les 5 autres clusters sont explicitement marqués comme non concluants par
    manque d'effectif.
  \end{block}
\end{frame}

% =============================================================================
% SLIDE 6 — Limites résiduelles et travaux futurs
% =============================================================================
\begin{frame}{Slide 6 --- Limites résiduelles + travaux futurs}
  \framesubtitle{8 limites identifiées par audit méthodologique post-Phase C}

  \begin{block}{Limites résiduelles (\texttt{docs/limitations.md} L10--L17)}
    \scriptsize
    \begin{tabular}{lll}
      \toprule
      \textbf{ID} & \textbf{Limite} & \textbf{Travail futur} \\
      \midrule
      L10 & H1 sil = 0.467 sur v4 (surentraînement siamois) & Hard-negative mining + curriculum \\
      L11 & Latence end-to-end & \textcolor{ewatgreen}{\checkmark{} Résolu (C-3, p95 = 13 ms)} \\
      L12 & OpenMax Unknown AUROC = 0.55 < 0.7 & Mahalanobis-OOD / Energy / ODIN \\
      L13 & Service graph $N = 6$ vs production $100+$ & Benchmark $N \geq 20$ \\
      L14 & Cross-cluster non validé & Stratégie B fine-tuning RCAEval (C-2) \\
      L15 & Pas de retraining cycle opérationnel & Drift d'embeddings → trigger retrain \\
      L16 & 17 features hardcodées & mRMR / SHAP-based pruning \\
      L17 & Ontologie sans audit SRE & 1 séance audit + intégration runbook \\
      \bottomrule
    \end{tabular}
  \end{block}

  \begin{block}{Stratégie globale}
    Chaque limite est documentée transparente \emph{et} accompagnée d'une
    proposition de fix technique. \textbf{Aucune limite cachée}, aucun fix
    suspect. Le rapport positionne EWAT comme une \textbf{contribution mature
    avec un programme de recherche clair} pour la suite (L13--L17 : 6+ mois
    de travaux futurs identifiés).
  \end{block}
\end{frame}

\end{document}
